\section{Manual Screening}
In smaller organisations, early-stage screening is still likely to be manually conducted, 
due to the cost of implementing and maintaining an automated screening system.
Recruiters and hiring managers will usually review a candidate's CV, along 
with application data, to compare against the job description to determine 
whether the applicant meets the minimum experience and skill requirements.
The screening stage is designed not to find the best candidate, 
but to instead to reject candidates that are considered unsuitable for the role, 
allowing for a more focused and efficient interview process.
In order to manage large applicant pools, recruiters may also introduce 
implicit filtering criteria in an attempt to simply reduce the number of 
candidates progressing into more resource-intensive interview stages.
While manual screening can be effective in recruitment, it often suffers 
from a range of limitations and biases that can reduce the overall quality, 
fairness and consistency of the process.

A primary issue revolves around biases due to time pressure and superficial 
evaluation, as recruiters usually spend little time on a first-pass, leading 
to them overlooking critical information.
The issues are further exacerbated by information overload and cognitive load 
bias, where excessive information leads to increased decision difficulties as 
well as a reliance on cognitive shortcuts rather than careful evaluation, as shown in research on information overload \cite{che2019}.
Furthermore, the repetitive task and fatigue effect suggests that when recruiters 
process multiple candidates in a session, their decision quality degrades.
This is supported by a PNAS study from Danziger et al., which highlights that favourable rulings drop 
sharply as decision-makers process more cases without breaks \cite{danziger2011}. 
Note that this work studied judicial parole decisions rather than hiring, but the same fatigue pattern may also apply to high-volume CV screening.
The mood and emotional state bias causes further issues as the Affect Infusion 
Model indicates that a recruiter's emotional state under stress can lead to a change 
in their perception or judgement of a candidate's suitability for that role \cite{forgas1995}.

Beyond these operational constraints, the screening process is also influenced by 
other distortions that introduce further bias.
The determinism of the recruitment process is affected by the contrast effect, as 
candidate evaluation may be unfairly affected by the quality of candidates that precede 
them \cite{palmer2014}.
This issue is made worse with the anchoring bias, where they over-rely on the first 
piece of information encountered, such as a previous employer or job title, which then 
disproportionately influences the rest of the screening \cite{buijsrogge2021}.
Halo and horn effects also tie in closely to this, where a recruiter generalises from 
one positive or negative attribute, such as a prestigious company, or a single typo, to 
the entire candidate profile \cite{thorndike1920}.
The previously mentioned issues are further worsened due to imperfect and exaggerated 
self-reported information where candidates may embellish or misrepresent their skills and 
experiences to appear more favourable \cite{dunning2004}.

Finally, manual screening is also vulnerable to systemic discrimination and group-based preferences.
For instance, field experiments by Bertrand and Mullainathan found that name 
and ethnic discrimination plays a key issue.
This is shown through their findings where “white sounding” names received around 
50\% more callbacks compared to “African American sounding” names despite identical qualifications \cite{bertrand2004}.
The main cause of this often comes down to similarity bias and in-group bias. This is the idea that
individuals tend to favour applicants who resemble them, often in terms of factors such as background, nationality
and educational route \cite{tajfel1986}.


The core limitations, as shown above, come down to time pressure, cognitive load, and applicant volume. 
The biases mentioned uncover the unfairness behind manual screening, more specifically the lack of consistency
and objectivity in this process.
As more candidates enter the application process, these psychological and procedural 
limitations compound further, leading to inconsistent outcomes, reduced fairness and 
also the potential for a poor selection quality.
Such challenges highlight the need for a more reliable, systematic and transparent 
screening method, more specifically automated, data-driven systems that are able to 
evaluate a candidate against a job requirement with better consistency, scalability, 
and a reduced vulnerability towards any potential human bias.
